\documentclass[11pt,a4paper]{article}
\usepackage[T1]{fontenc}
\usepackage[utf8]{inputenc}
\usepackage[english]{babel}
\usepackage{lmodern}
\usepackage{amsmath,amssymb}
\usepackage{xcolor}
\usepackage[margin=2.5cm]{geometry}
\usepackage{microtype}
\definecolor{PFFBlue}{HTML}{203E6A}
\setlength{\parindent}{0pt}
\setlength{\parskip}{0.45em}
\begin{document}
\section*{\textcolor{PFFBlue}{Lemma 5.1 --- Zero Integral Implies Zero Almost Everywhere}}

\textbf{Statement.} Let $(S,\Sigma,\mu)$ be a measure space and let
$f\in(m\Sigma)^+$ satisfy $\mu(f)=0$. Then
\[
\mu(\{f>0\})=\mu(\{s\in S:f(s)>0\})=0.
\]

\textbf{Proof.}\par
We argue by contradiction. Suppose that $\mu(\{f>0\})>0$ and define
\[
A_n:=\left\{f>\frac1n\right\}\in\Sigma,\qquad n\in\mathbb N.
\]
The sets are measurable because $f$ is $\Sigma$-measurable. As we will see in
HW~46,
\[
\{f>0\}
=\bigcup_{n=1}^{\infty}\left\{f>\frac1n\right\}
=\lim_{n\to\infty}\left\{f>\frac1n\right\}.
\]
Moreover, $A_n\subseteq A_{n+1}$ because $1/(n+1)<1/n$. By continuity
from below (Lemma~1.5),
\[
0<\mu(\{f>0\})=\lim_{n\to\infty}\mu(A_n).
\]
Since $(\mu(A_n))_{n\in\mathbb N}$ is nondecreasing and its limit is
strictly positive, there exists $\bar n\in\mathbb N$ such that
$\mu(A_n)>0$ for every $n\ge\bar n$. Fix $n=\bar n$ and consider the
nonnegative simple function
\[
h:=\frac1{\bar n}\mathbf1_{A_{\bar n}}\in SF^+.
\]
We have $h\le f$: if $s\in A_{\bar n}$, then
$h(s)=1/\bar n<f(s)$; otherwise $h(s)=0\le f(s)$.
Therefore, by the definition of the positive Lebesgue integral,
\[
\begin{aligned}
0=\mu(f)
&=\sup\{\mu_0(g):g\in SF^+,\ g\le f\}\\
&\ge\mu_0(h)
=\frac1{\bar n}\mu(A_{\bar n})>0.
\end{aligned}
\]
The inequality holds because $h$ is an admissible simple function; the
following equality uses the definition of $\mu_0$, and the last inequality
follows from $\mu(A_{\bar n})>0$. This is a contradiction. Thus
$\mu(\{f>0\})=0$, i.e., $f=0$ $\mu$-a.e.

Therefore, the statement is proved. $\square$
\end{document}
